\section{Mean Removal and What the Objective Actually Learns}
\label{sec:mean_removal_from_loss}

This section answers a foundational question: if the token-mean of activations is a random variable that depends on context, why is it nevertheless plausible that removing (or separating) the mean does not destroy the learning signal?
We proceed from the training objective of LLMs, derive the exact gradient expressions, and establish which parts of the signal correspond to the mean (first moment) and which correspond to token-wise variability (second-order structure).

\subsection{Training Objective and Notation}

Large language models are trained by minimizing next-token cross-entropy.
Let $\mathcal{D}$ denote the data distribution over sequences, and let the model define conditional probabilities
\begin{equation}
p_\theta(w_{t+1}\mid w_{\le t}) = \mathrm{softmax}(z_t)_ {w_{t+1}},
\end{equation}
where $z_t\in\mathbb{R}^V$ are logits over vocabulary size $V$.
The population loss is
\begin{equation}
\label{eq:ce_loss}
\mathcal{L}(\theta)=\mathbb{E}_{(w_1,\dots,w_T)\sim\mathcal{D}}\Bigg[-\sum_{t=1}^{T-1}\log \mathrm{softmax}(z_t)_{w_{t+1}}\Bigg].
\end{equation}

Fix a layer $\ell$ and focus on a single linear projection within that layer (the argument applies to attention and MLP projections alike).
Let $x_t\in\mathbb{R}^H$ be the input activation at token $t$ to this projection, and define the projection
\begin{equation}
\label{eq:linear_proj}
y_t = W x_t + b,\qquad W\in\mathbb{R}^{d\times H},\; b\in\mathbb{R}^d.
\end{equation}
Downstream nonlinearities and further layers map $\{y_t\}$ to logits $\{z_t\}$.
Let $\delta_t\in\mathbb{R}^d$ denote the backpropagated gradient at this projection:
\begin{equation}
\label{eq:delta_def}
\delta_t \;\triangleq\; \frac{\partial \mathcal{L}}{\partial y_t}.
\end{equation}
Then the per-token gradient of $\mathcal{L}$ with respect to $W$ is exactly
\begin{equation}
\label{eq:grad_outer}
\nabla_W \mathcal{L} \;=\; \sum_{t} \delta_t x_t^\top,
\qquad
\nabla_b \mathcal{L} \;=\; \sum_t \delta_t.
\end{equation}
This is the precise ``outer-product'' structure used throughout optimization.

\subsection{Mean Decomposition and Exact Gradient Split}

Let $N$ denote the number of tokens contributing to the gradient in a batch (or in an expectation over data).
Define the \emph{token-mean activation} and \emph{centered activation}:
\begin{equation}
\label{eq:mean_center_def}
\mu_x \;\triangleq\; \frac{1}{N}\sum_{t=1}^N x_t,
\qquad
\tilde{x}_t \;\triangleq\; x_t - \mu_x,
\qquad
\sum_{t=1}^N \tilde{x}_t = 0.
\end{equation}
Substitute $x_t=\mu_x+\tilde{x}_t$ into \eqref{eq:grad_outer}. This yields an \emph{exact identity}:
\begin{equation}
\label{eq:grad_split}
\nabla_W \mathcal{L}
=
\sum_t \delta_t(\mu_x+\tilde{x}_t)^\top
=
\Big(\sum_t \delta_t\Big)\mu_x^\top
+
\sum_t \delta_t \tilde{x}_t^\top.
\end{equation}
Define the mean of backprop signals
\begin{equation}
\label{eq:mu_delta_def}
\mu_\delta \;\triangleq\; \frac{1}{N}\sum_{t=1}^N \delta_t.
\end{equation}
Then \eqref{eq:grad_split} can be written as
\begin{equation}
\label{eq:grad_split_compact}
\frac{1}{N}\nabla_W \mathcal{L}
=
\underbrace{\mu_\delta \mu_x^\top}_{\text{rank-1 mean--mean term}}
+
\underbrace{\frac{1}{N}\sum_t \delta_t\tilde{x}_t^\top}_{\text{cross-covariance term}}.
\end{equation}

\paragraph{Interpretation.}
Equation~\eqref{eq:grad_split_compact} is the end-to-end decomposition that matters for learning.
The first term is rank-one and depends only on the \emph{means} $(\mu_\delta,\mu_x)$.
The second term depends on token-wise \emph{variations} $\tilde{x}_t$ and is a cross-covariance between
backprop signals and token-specific deviations.
If the optimization objective requires learning fine-grained distinctions between tokens and contexts, that information must appear in the cross-covariance term.

\subsection{A Fully Rigorous Invariance Statement: Centering with Bias Compensation}

Mean removal in the forward pass changes the function \emph{unless} we compensate for the removed mean through an equivalent parameter transformation.
We now state and prove the exact equivalence.

\begin{theorem}[Exact functional equivalence under mean-centering with bias update]
\label{thm:centering_equiv}
Fix a batch (or context) producing activations $\{x_t\}_{t=1}^N$ and mean $\mu_x$.
Define centered activations $\tilde{x}_t = x_t - \mu_x$.
Consider the linear map \eqref{eq:linear_proj}.
Define an adjusted bias
\begin{equation}
\label{eq:bprime}
b' \;\triangleq\; b + W\mu_x.
\end{equation}
Then for every token $t$,
\begin{equation}
\label{eq:exact_equiv}
W x_t + b \;=\; W\tilde{x}_t + b'.
\end{equation}
Consequently, replacing $x_t$ by $\tilde{x}_t$ \emph{and} replacing $b$ by $b'$ leaves all downstream activations and logits unchanged, hence leaves the cross-entropy loss \eqref{eq:ce_loss} unchanged.
\end{theorem}

\begin{proof}
By definition, $x_t=\tilde{x}_t+\mu_x$. Substitute into the left-hand side:
\[
W x_t + b = W(\tilde{x}_t+\mu_x)+b = W\tilde{x}_t + (W\mu_x+b).
\]
With $b'$ defined by \eqref{eq:bprime}, we obtain \eqref{eq:exact_equiv}.
Since $y_t$ is unchanged for all $t$, every downstream deterministic computation produces identical logits $z_t$, and therefore the softmax probabilities and the loss \eqref{eq:ce_loss} are identical.
\end{proof}

\paragraph{Consequence.}
Theorem~\ref{thm:centering_equiv} is the precise sense in which ``mean removal need not destroy predictive information'':
it is a \emph{reparameterization} of the same function, provided the mean is forwarded (absorbed) as an equivalent bias contribution at the next linear boundary.
This theorem is exact and does not assume anything about the distribution of $\mu_x$.

\subsection{What If We Center Without Compensation?}

If we replace $x_t$ by $\tilde{x}_t$ but keep the original bias $b$, then the linear output becomes
\[
y_t' = W\tilde{x}_t + b = y_t - W\mu_x,
\]
which is a constant shift in $y_t$ across tokens for the given batch.
This shift generally changes logits and therefore changes the loss.
Thus ``centering is free'' is not true without qualification.
The correct statement is: centering is information-preserving conditional on forwarding/absorbing the mean into a separate channel (e.g., adjusted bias, residual carrier, or explicit structure stream).

\subsection{Why ``Useful Information Lives in Variability''}

Equation~\eqref{eq:grad_split_compact} now makes the key point precise.
The only part of the gradient that can encode token- and context-specific distinctions is the cross-covariance term:
\begin{equation}
\label{eq:cross_cov}
C_{\delta x}\;\triangleq\;\frac{1}{N}\sum_{t=1}^N \delta_t\tilde{x}_t^\top.
\end{equation}
The rank-one mean--mean term $\mu_\delta\mu_x^\top$ is shared across tokens and cannot express within-batch variation.

In particular, if $\mu_\delta$ is small (noise-level)---as observed empirically for gradient outputs in our measurements---then the mean--mean term is negligible and learning is driven almost entirely by $C_{\delta x}$.
Even when $\mu_\delta$ is not identically zero, the rank-one term represents a global shared update direction, whereas $C_{\delta x}$ captures structured correlations between errors and token-level deviations, which is the mechanism required to learn conditional dependencies beyond base rates.

\subsection{Summary}

Starting from the cross-entropy loss \eqref{eq:ce_loss}, we derived an exact gradient decomposition \eqref{eq:grad_split_compact} showing that updates split into a rank-one mean--mean component and a cross-covariance component.
We then proved an exact invariance theorem (Theorem~\ref{thm:centering_equiv}) showing that mean-centering can be made function-preserving when the removed mean is forwarded as an equivalent bias at the next linear boundary.
Therefore, mean separation is not inherently information-destroying; information is lost only if the mean channel is discarded rather than separated.